% ============================================================================
% Chapter I -- Preliminary Study
% ============================================================================

% ----------------------------------------------------------------------------
% Introduction G\'en\'erale (unnumbered, before Chapter I)
% ----------------------------------------------------------------------------
\chapter*{General Introduction}
\addcontentsline{toc}{chapter}{General Introduction}
\markboth{General Introduction}{}

The digital era has brought unprecedented connectivity---and with it, an
equally unprecedented exposure of sensitive data. According to industry
estimates, over 26~billion records were exposed in data breaches during 2023
alone, and the trend has only accelerated since~\cite{ibm-breach-report}. Every
month, new breach compilations, stealer logs, and combolists surface on
darknet forums, paste sites, and Telegram channels, making personally
identifiable information (PII) and credentials available to threat actors at
an industrial scale.

In response, a discipline known as \textbf{Open-Source Intelligence (OSINT)}
has emerged as a critical capability for cybersecurity defenders. OSINT
platforms allow analysts to search leaked data proactively---not to exploit it,
but to assess exposure, notify affected parties, and harden defenses before
attackers strike. Commercial platforms such as IntelligenceX, DeHashed, and
Snusbase provide this capability, but they are often expensive, opaque in
their architecture, and inaccessible to students and small security teams.

It is in this context that \leaklens{} was conceived: an OSINT intelligence
platform that enables authorized security researchers to search breached
databases by email, username, IP address, password, name, hash, or domain.
Results include credential exposure records enriched with supplementary
intelligence from public APIs. The project also demonstrates defensive
software engineering---hardening the application itself against common web
attacks.

This report documents the full engineering cycle, with emphasis on:

\begin{itemize}[nosep]
  \item The cybersecurity context and need for accessible OSINT tooling;
  \item Functional and non-functional requirements;
  \item Security-oriented architecture (JWT, RBAC, rate limiting, CSP);
  \item Production realization on Vercel and a VPS.
\end{itemize}

The report is organized in four chapters. \chapref{chap:preliminary} presents
the preliminary study, analyzes existing solutions, and justifies the choices
made. \chapref{chap:requirements} details the functional requirements and
technology stack. \chapref{chap:architecture}, the core of this report,
describes the security-oriented architecture design.
Finally, \chapref{chap:implementation} demonstrates the concrete realization
through screenshots and production metrics.

\clearpage

% ============================================================================
\chapter{Preliminary Study}
\label{chap:preliminary}
% ============================================================================

\section*{Introduction}

Before designing any technical solution, it is essential to understand the
context in which it operates, to analyze the tools already available on the
market, and to identify the gaps that the project aims to fill. This first
chapter establishes the foundations of the project: its context, problem
statement, and value for cybersecurity research.

% ----------------------------------------------------------------------------
\section{General Context}
\label{sec:context}
% ----------------------------------------------------------------------------

\subsection{The Cybersecurity Landscape (2024--2026)}

Between 2024 and 2026, the volume and \emph{reuse} of leaked breach data
became a defining feature of the threat landscape. Incidents no longer end
when a company publishes a breach notification---the stolen records live on
independently, copied into combolists, stealer-log archives, and searchable
databases traded on criminal forums, Telegram channels, and paste sites.

Several trends illustrate the scale of the problem:

\begin{itemize}[nosep]
  \item \textbf{2024 --- aggregation at industrial scale:} The ``Mother of All
        Breaches'' (MOAB) consolidated an estimated 26~billion records from
        dozens of prior incidents into a single searchable corpus, proving
        that breach data is permanently recycled rather than
        contained~\cite{moab}.
  \item \textbf{2025 --- stealer logs as a primary source:} Infostealer
        malware families (RedLine, Raccoon, Vidar, Lumma) harvest
        browser-saved credentials, session cookies, and autofill data from
        infected machines. These \emph{stealer logs} are sold in bulk and
        often contain live session tokens, making them more dangerous than
        static password dumps from older database breaches.
  \item \textbf{2026 --- sustained corporate and public-sector impact:}
        Healthcare providers, financial institutions, and government agencies
        continue to report large-scale exposures. IBM estimates the global
        average cost of a data breach at USD~4.88~million in
        2024~\cite{ibm-breach-report}, driven largely by post-breach
        remediation, regulatory fines, and customer churn.
\end{itemize}

\paragraph{What breach compilations contain.}
A typical leak dataset may include email addresses, usernames, plaintext or
reused passwords, password hashes (MD5, SHA-1, bcrypt), IP addresses, phone
numbers, physical addresses, and internal corporate identifiers. NIST
classifies much of this material as personally identifiable information
(PII)~\cite{nist-breach}. Once published, a single credential pair can appear
in dozens of derivative compilations within weeks.

\paragraph{How attackers exploit leaked data.}
Threat actors treat breach data as reusable ammunition rather than a one-time
event. Common abuse patterns include:

\begin{enumerate}[nosep]
  \item \textbf{Credential stuffing:} Automated bots test email/password
        pairs harvested from one breach against banking portals, SaaS
        platforms, and corporate VPN gateways. Because users frequently reuse
        passwords, a leak from a low-value forum can unlock high-value
        corporate accounts.
  \item \textbf{Account takeover (ATO):} Valid credentials enable direct
        access to email, cloud storage, and social media. Attackers pivot
        from ATO to password-reset flows on linked services, expanding the
        blast radius without additional malware.
  \item \textbf{Targeted phishing and social engineering:} Knowing that a
        victim's email appeared in a specific breach (e.g., a healthcare or
        payroll provider) allows attackers to craft convincing lures
        referencing real incident details, dramatically increasing click and
        compromise rates.
  \item \textbf{Initial access and ransomware:} Breach-derived credentials
        are sold on initial-access broker markets. Ransomware groups purchase
        valid VPN or remote-desktop logins to bypass perimeter defenses
        entirely.
  \item \textbf{Identity fraud and financial crime:} Combined PII fields
        (name, address, phone, date of birth) support synthetic identity
        creation, SIM-swap preparation, and fraudulent loan or tax
        applications.
\end{enumerate}

\paragraph{Why the danger persists.}
Even when organizations rotate passwords after an incident, historical breach
records remain valuable: weak hashes can be cracked offline, reused passwords
can be tested indefinitely, and stealer logs may capture credentials \emph{after}
a rotation. Defenders therefore need continuous visibility into whether their
assets---employee emails, corporate domains, public IP ranges---appear in
newly surfaced leak compilations, not merely whether a breach was announced
years ago.

\subsection{Problem Statement}
\label{sec:problematic}

The core problem is asymmetric: attackers can acquire and weaponize breach
data cheaply and at scale, while defenders often lack timely, actionable
visibility into the same datasets. A security team may learn that \emph{a}
breach occurred, yet remain unable to answer precise questions that determine
real risk:

\begin{itemize}[nosep]
  \item Do our employees' corporate email addresses appear with plaintext
        passwords in stealer logs or combolists circulating this month?
  \item Has a contractor's reused password from an unrelated consumer breach
        exposed our VPN or Microsoft~365 tenant to credential stuffing?
  \item Does a newly registered phishing domain overlap with subdomains or
        WHOIS data tied to our organization's prior exposures?
\end{itemize}

Without answers, organizations operate blind during the critical window between
data exfiltration and public disclosure---often months or years---during which
threat actors actively monetize the stolen records.

\paragraph{Defender use cases.}
Proactive breach-data search supports three operational needs:

\begin{enumerate}[nosep]
  \item \textbf{Exposure assessment:} Map which identities, domains, and IP
        addresses within an organization's scope appear in known leak
        compilations, prioritizing remediation (forced resets, MFA rollout,
        session revocation) before attackers exploit the overlap.
  \item \textbf{Incident response:} During or after a suspected breach,
        rapidly determine whether specific assets, hashes, or internal
        identifiers have surfaced in external datasets, shortening
        investigation timelines.
  \item \textbf{Threat intelligence enrichment:} Correlate a suspicious email,
        domain, or IP with historical leak context to assess whether an
        observed login attempt likely originates from a recycled credential
        rather than a novel attack vector.
\end{enumerate}

\paragraph{Gap in existing tooling.}
Current market offerings leave significant blind spots:

\begin{itemize}[nosep]
  \item \textbf{Notification without substance:} Services such as Have I Been
        Pwned~\cite{hibp} confirm that an email appeared in a named breach
        but withhold the underlying credential data and enrichment context
        needed for thorough triage.
  \item \textbf{Prohibitive cost:} Professional engines like
        IntelligenceX~\cite{intelx} and DeHashed~\cite{dehashed} return
        full records but charge subscription fees that exclude students,
        independent researchers, and small security teams.
  \item \textbf{Opaque engineering:} Commercial platforms do not expose their
        aggregation, validation, or access-control design, limiting their
        value as reference implementations for academic or audit purposes.
  \item \textbf{No integrated, self-hosted option:} Few tools combine
        multi-type breach search, supplementary OSINT enrichment, role-based
        access control, and containerized deployment in a single transparent
        stack.
\end{itemize}

\leaklens{} targets this gap directly: an OSINT platform that lets
authorized analysts query real breach records across seven search types,
enriched with public API data, while demonstrating how such a sensitive
capability should be engineered---invite-gated registration, JWT-based
authorization, rate limiting, and no local storage of raw credentials beyond
audit metadata.

% ----------------------------------------------------------------------------
\section{Study of Existing Solutions}
\label{sec:existing}
% ----------------------------------------------------------------------------

\subsection{Solution~1: Have I Been Pwned (HIBP)}

Have I Been Pwned, created by security researcher Troy Hunt, is one of the
most widely used breach-notification services in the world. Users enter an
email address and HIBP reports which known breaches include that
address~\cite{hibp}.

\textbf{Strengths:}
\begin{itemize}[nosep]
  \item Free and publicly accessible;
  \item Covers over 700 breaches and 13+ billion accounts;
  \item Offers a \texttt{Pwned Passwords} API for password checking;
  \item Respected and trusted in the security community.
\end{itemize}

\textbf{Limitations:}
\begin{itemize}[nosep]
  \item Only reveals breach \emph{names}---never actual credentials;
  \item Search limited to email and phone; no username, IP, or domain search;
  \item No enrichment (subdomains, WHOIS, reputation);
  \item API is rate-limited and the paid tier is enterprise-priced.
\end{itemize}

\subsection{Solution~2: IntelligenceX (IntelX)}

IntelligenceX is a professional OSINT search engine that indexes data from
the darknet, paste sites, public records, and breach compilations. It
supports searching by email, domain, IP, Bitcoin address, and
more~\cite{intelx}.

\textbf{Strengths:}
\begin{itemize}[nosep]
  \item Extremely powerful: searches darknet, Tor, I2P, paste sites;
  \item Returns actual file contents and credential data;
  \item Supports advanced search operators and historical data;
  \item Professional-grade API with comprehensive documentation.
\end{itemize}

\textbf{Limitations:}
\begin{itemize}[nosep]
  \item Expensive: professional plans start at \$2,000+/year;
  \item Free tier is heavily restricted (limited searches per day);
  \item Closed-source, proprietary architecture;
  \item Enterprise-focused; inaccessible for academic projects.
\end{itemize}

\subsection{Solution~3: DeHashed}

DeHashed is a subscription-based breach search engine that allows users to
search leaked databases by email, username, IP address, name, phone, and
more~\cite{dehashed}.

\textbf{Strengths:}
\begin{itemize}[nosep]
  \item Affordable entry point (\$5.49/month for basic plan);
  \item Returns actual credentials (passwords, hashes);
  \item API access included in paid plans;
  \item Supports multiple search types (email, username, IP, name, phone).
\end{itemize}

\textbf{Limitations:}
\begin{itemize}[nosep]
  \item No free tier; requires subscription for any meaningful use;
  \item No enrichment capabilities (WHOIS, subdomains, reputation);
  \item No self-hosted deployment option;
  \item No role-based access control or invite-gated registration.
\end{itemize}

% ----------------------------------------------------------------------------
\section{Comparative Study}
\label{sec:comparative}
% ----------------------------------------------------------------------------

\tabref{tab:comparison} compares the existing solutions with \leaklens{}
across several key criteria.

\begin{table}[ht]
\centering
\caption{Comparative table of OSINT breach-search platforms.}
\label{tab:comparison}
\small
\begin{tabularx}{\textwidth}{@{} L{3.8cm} C{2cm} C{2cm} C{2cm} C{2cm} @{}}
\toprule
\textbf{Criterion} & \textbf{HIBP} & \textbf{IntelX} & \textbf{DeHashed}
  & \textbf{LeakLens} \\
\midrule
Free access                 & \cmark & $\sim$ & \xmark & \cmark \\
Credential data returned    & \xmark & \cmark & \cmark & \cmark \\
Multiple search types       & \xmark & \cmark & \cmark & \cmark \\
Wildcard search             & \xmark & \cmark & \cmark & \cmark \\
Enrichment (WHOIS, subs.)   & \xmark & $\sim$ & \xmark & \cmark \\
Self-hosted / Docker        & \xmark & \xmark & \xmark & \cmark \\
RBAC \& invite-gated signup & \xmark & \xmark & \xmark & \cmark \\
Security-hardened backend   & N/A    & N/A    & N/A    & \cmark \\
Open architecture           & \xmark & \xmark & \xmark & \cmark \\
Affordable for students     & \cmark & \xmark & $\sim$ & \cmark \\
\bottomrule
\end{tabularx}

\vspace{4pt}
\footnotesize \cmark\ = Yes \quad $\sim$ = Partial \quad \xmark\ = No
\end{table}

This analysis highlights that, unlike commercial solutions, the proposed
system offers full control over the stack and security model---its main
advantage for academic and professional use.

% ----------------------------------------------------------------------------
\section{Proposed Solution: LeakLens}
\label{sec:proposed}
% ----------------------------------------------------------------------------

\subsection{Project Vision}

\leaklens{} is an OSINT platform designed as a complete demonstrator of
modern, security-first software engineering.

The vision of \leaklens{} rests on three fundamental pillars:

\begin{enumerate}[nosep]
  \item \textbf{Intelligence:} A real, functional search platform covering
        seven search types across billions of breached records, enriched with
        supplementary data from free public APIs.
  \item \textbf{Security:} A defense-in-depth architecture implementing
        industry best practices---JWT, RBAC, rate limiting, CSP, HSTS,
        constant-time authentication, and Docker hardening.
  \item \textbf{Simplicity:} A lean, containerized deployment that can be
        reproduced on any VPS with Docker Compose, paired with a
        Vercel-hosted frontend for zero-maintenance delivery.
\end{enumerate}

\subsection{Project Objectives}

\begin{itemize}[nosep]
  \item Develop a full-stack OSINT platform with Next.js (frontend) and
        FastAPI (backend);
  \item Integrate Snusbase as the primary leak-data engine, supplemented
        by free enrichment APIs;
  \item Implement comprehensive authentication and authorization (JWT,
        RBAC, invite-code gating);
  \item Harden the application against common web attacks (enumeration,
        brute force, XSS, CSRF);
  \item Containerize with Docker using Alpine multi-stage builds, non-root
        execution, and vulnerability scanning;
  \item Deploy in production on a VPS (backend) and Vercel (frontend) with
        automatic TLS.
\end{itemize}

\subsection{High-Level Architecture}

\figref{fig:high-level-arch} presents the high-level architecture. The
frontend runs on Vercel; the backend and database run on a VPS inside Docker
Compose. A Next.js rewrite proxy routes API calls server-side, avoiding
mixed-content issues when the backend has no custom HTTPS domain.

\begin{figure}[ht]
\centering
\tikzfig{%
\begin{tikzpicture}[
  node distance=1.1cm and 1.6cm,
  box/.style={draw,rounded corners,minimum width=2.4cm,minimum height=0.95cm,
    align=center,font=\footnotesize},
  arrow/.style={-{Stealth[length=2.5mm]},thick},
]
  \node[box,fill=blue!10] (browser) {Browser};
  \node[box,fill=green!10,right=of browser] (vercel)
    {Vercel\\Next.js};
  \node[box,fill=orange!10,right=of vercel] (vps)
    {VPS\\FastAPI + PG};
  \node[box,fill=red!10,right=of vps] (snusbase) {Snusbase};
  \node[box,fill=yellow!10,below=0.6cm of snusbase] (freeapis)
    {Enrichment APIs};

  \draw[arrow] (browser) -- node[above,font=\tiny]{HTTPS} (vercel);
  \draw[arrow] (vercel) -- node[above,font=\tiny]{Proxy} (vps);
  \draw[arrow] (vps) -- node[above,font=\tiny]{HTTPS} (snusbase);
  \draw[arrow] (vps) |- (freeapis);
\end{tikzpicture}%
}
\caption{High-level system architecture.}
\label{fig:high-level-arch}
\end{figure}

% ----------------------------------------------------------------------------
\section*{Conclusion}

This first chapter established the project context, analyzed existing
OSINT platforms, and positioned the proposed solution as a practical,
security-conscious alternative for students and small security teams. The
next chapter details requirements and technology choices.
